\subsection{A nonzero coherent-norm obstruction with every finite norm test zero}

\begin{statement}{A.181}{Proved}
{An evaluated unit support class and a locally trivial one-variable obstruction}
Retain the fixed tower of A.179, put
\[
 \begin{gathered}
 k=K_1=\Q(\alpha),\quad \alpha^3=3,\quad
 d=649,\quad A=\rho_1=d^8\alpha,\\
 U=\alpha^2-2,\quad B=U^3,
 \quad K_m=k(A^{1/3^{m-1}}).
 \end{gathered}
\tag{A.181.1}
\]
All roots and embeddings are the positive real ones already fixed.
Write \(\mathcal U_m=\mathcal O_{K_m}^{\times}\),
\(\mathcal U_{m,S}=\mathcal O_{K_m,S_m}^{\times}\), with
\(S=\{2,3,11,59,\infty\}\), and let a subscript three on a
hat mean pro-three completion.  The exact completed global-unit norm
images in \(\widehat{\mathcal U}_{1,3}=U^{\Z_3}\) are
\[
 N_{K_2/k}\widehat{\mathcal U}_{2,3}=U^{3\Z_3},\qquad
 N_{K_3/k}\widehat{\mathcal U}_{3,3}=U^{9\Z_3}.
\tag{A.181.2}
\]
The actual element \(B\) gives an element of exact order three in
each of the following two kernels, with its indicated image under
completion:
\[
\begin{aligned}
 &\ker\left(\mathcal U_{1,S}/N_{K_3/k}\mathcal U_{3,S}
       \longrightarrow k^\times/N_{K_3/k}K_3^\times\right),\\
 &\ker\left(\widehat{\mathcal U}_{1,S,3}/
                    N_{K_3/k}\widehat{\mathcal U}_{3,S,3}
       \longrightarrow \widehat{k^\times}_3/
                    N_{K_3/k}\widehat{K_3^\times}_3\right).
\end{aligned}
\tag{A.181.3}
\]
For every finite \(m\ge1\), \(B\) is an actual field norm from
\(K_m\) to \(k\).  At every place of \(k\) it belongs to a
compatible local norm family through this entire tower.  Nevertheless
both global inverse fibres are empty:
\[
\begin{aligned}
 \{(b_m)\in\varprojlim_N K_m^\times:b_1=B\}&=\varnothing,\\
 \{(b_m)\in\varprojlim_N\widehat{K_m^\times}_3:b_1=[B]\}
       &=\varnothing.
\end{aligned}
\tag{A.181.4}
\]

For the geometric one-variable quotient \(P\) of A.77, now over
\(k\) and at coefficient three, the absolute Kummer pair
\((\kappa_3(A),\kappa_3(B))\) therefore has a nonzero continuous
obstruction
\[
 \mathfrak o=\delta\kappa_3(B)
      \in H^2_{\mathrm{cts}}(G_k,D_A).
\tag{A.181.5}
\]
Its restriction to every completion of \(k\) is zero, and its
image in every individual root-torsor quotient is zero.  The proof
specifies \(D_A\), the exact coefficient sequence, the actual
quotient maps, and the nonzero norm-kernel derived-limit class which
maps to \(-\mathfrak o\).  These assertions concern the displayed
Kummer pair over \(k\), with its specified norm maps.
\end{statement}

\begin{proof}
\emph{The evaluated global-unit and fixed-support norm images.}
The inverse \(U^{-1}=4+3\alpha+2\alpha^2\) and the identity
\(N_{k/\Q}U=1\) were checked by literal field arithmetic in
A.159.  The unit \(U\) is a pro-three basis of the rank-one
unit group of \(k\), by the residue and saturation argument in
A.165.  The four units of A.159 give a pro-three basis of the
global units of \(K_2\).  Their relative norm exponents and their
normalized local characters are
\[
\begin{gathered}
 N_{K_2/k}u_i=U^{k_i},\qquad
 (k_i)=(3,57,-3,1614),\qquad
 g_i=k_i/3=(1,19,-1,538),\\
 \lambda_q(u_iU^{-g_i})=1\quad(q=11,59),\qquad
 \lambda_q(U)=\lambda_q(\alpha)^2\ne1.
\end{gathered}
\tag{A.181.6}
\]
Here \(\lambda_q\) is the cubic character of the residue field
\(\F_{q^2}\) at the quadratic branch, and
\(\lambda_q(\alpha)\) has order three.  The two polynomial
residue equations are
\(\alpha^2-2\alpha+4=0\) at eleven and
\(\alpha^2-10\alpha+41=0\) at fifty-nine.
All identities in A.181.6 are the certified actual maps
A.159.3--8, with their characters expressed in these generators.

For \(v=\prod_i u_i^{z_i}\), \(z_i\in\Z_3\), put
\(m=\sum_i g_i z_i\).  Then
\[
 Nv=U^{3m},\qquad
 \lambda_{11}(v)=\lambda_{11}(U)^m,\qquad
 \lambda_{59}(v)=\lambda_{59}(U)^m.
\tag{A.181.7}
\]
Since \(g_1=1\), the first equality in A.181.2 follows.
Each changed-prime completion of \(K_3/K_2\) is a single
totally ramified cubic extension.  A norm of a local unit has cube
residue.  A completed global unit which is such a norm therefore
has \(m\in3\Z_3\) by A.181.7.  Norm transitivity proves the
inclusion of the second image in \(U^{9\Z_3}\).
The opposite inclusion has the explicit scalar preimage: the
element \(U\) viewed in \(K_3\) has norm \(U^9\).
This proves A.181.2 exactly.

The fixed-support obstruction uses the larger group, so we retain
its kernel map.  Suppose a completed \(S\)-unit \(z\) of
\(K_3\) had norm \([U^3]\) to \(k\).  Put
\(v=N_{K_3/K_2}z\).  The scalar \(U\in K_2\) has relative
norm \(U^3\), so \(v/U\) lies in the completed relative
norm-one \(S\)-unit group of A.159.12.  That group has valuation
zero and trivial cubic unit character at both changed-prime branches.
Thus \(v\) has the nonidentity characters of \(U\), although
it would be a local norm from the totally ramified cubic above it.
Its valuation there is zero, so any local norm preimage also has
valuation zero, and its residue must be a cube.  This is a contradiction.
It excludes an actual \(S\)-unit norm as well, by its natural image
in the completed norm diagram.  The cube \(B^3=U^9\) is the
literal norm of the scalar \(U\in\mathcal U_3\).
After the field-norm assertion below is proved, this establishes
both exact order-three claims in A.181.3.

\emph{Every finite field-norm map and every local inverse fibre.}
Fix a place \(v\) of \(k\) above \(q=11\) or \(59\).
The local valuation of \(A=d^8\alpha\) is eight.  For
\(e=3^{m-1}\), the polynomial \(X^e-A\) therefore gives a
single totally ramified field of degree \(e\).  At the quadratic
branch the residue field has order \(q^2\), with
\(v_3(q^2-1)=1\).  The residue norms of units are exactly
the \(e\)-th powers.  This last equality is explicit: a scalar
lift realizes a prescribed residue \(e\)-th power, and the
\(e\)-power map is an automorphism on the principal units because
\(q\nmid e\).  For every \(m\ge2\),
\[
 (\F_{q^2}^{\times})^{3^{m-1}}
       =(\F_{q^2}^{\times})^3,
       \qquad \overline B=\overline U^3
                  \in(\F_{q^2}^{\times})^3.
\tag{A.181.8}
\]
Thus \(B\) is a unit norm at every depth.  At the linear branch
the residue field is \(\F_q\), whose multiplicative order is
prime to three, so the same argument gives every unit as a norm.

At three the local field \(k_v\) has degree three over
\(\Q_3\), contains no \(\mu_3\), and has
\(v_{k_v}(A)=1\).  The actual local Artin/norm diagram of
A.140 identifies the norm image from each pure-root field with
the image of its root stabilizer in the abelianization of its
Galois closure.  More explicitly that group is
\(C_{3^{m-1}}\rtimes H\), with \(H\) the cyclotomic
root stabilizer.  Its element congruent to \(-1\) modulo three
makes the translation group a commutator subgroup.  Hence
\(H\to(C_{3^{m-1}}\rtimes H)^{\rm ab}\) is onto.
The local reciprocity norm map, followed by a norm from the
splitting field, consequently gives a pure-field norm preimage
of every target, by the explicit construction of A.140.8--9
and \cite[III, Theorem 3.5]{MilneCFT}.
All these local fields are totally ramified.  A preimage of the
unit \(B\) is therefore a unit.

At every finite place not over \(3,11,59\), \(A\) and
\(3\) are units.  The root algebras are unramified, since
\(X^{3^{m-1}}-A\) and its derivative define an etale integral
algebra.  The target \(B\) is a unit, and unit norms in each
unramified field factor are onto.  One can put the target in one
factor and one in all other factors.  At a real place use the
chosen real root factor; at a complex place choose a compatible
embedding of the root tower and put the target in that factor.
These are the actual product norm maps on the etale algebras.

At every finite place, unit groups and their finite products are
compact.  The preceding argument supplies a unit preimage at
every finite top height.  Its successive norms solve every
finite initial collection of compatibility equations.  The
closed norm equations in the product of these compact unit
groups therefore have the finite intersection property.
They give compatible actual local unit families, and hence
compatible completed local families.  The archimedean choices
were already compatible along the chosen embeddings.

Here is the exact global norm comparison.  Let
\[
 \Omega_m=\Q(\mu_{3^m},\rho_m),\quad
 G_m=(3\Z/3^m)\rtimes(\Z/3^m)^\times,
 \quad H_m=(\Z/3^m)^\times.
\tag{A.181.9}
\]
These are respectively \(\operatorname{Gal}(\Omega_m/k)\)
and \(\operatorname{Gal}(\Omega_m/K_m)\).
Indeed the full global group and the full local group at three
both have order \(3^m\varphi(3^m)\), as proved in
A.162.12: the Eisenstein cubic root of \(a=3d^{24}\)
cannot lie in an abelian local cyclotomic extension, so its
Kummer order over that field is \(3^m\).
The translations fixing \(\rho_1\) are exactly
\(3\Z/3^m\).  Thus the decomposition subgroup over the
unique place of \(k\) above three is all of \(G_m\).

For the norm-one torus \(T_m=R^1_{K_m/k}\mathbf G_m\),
the exact lattice and localization maps are
\[
\begin{gathered}
 X_m=\Z[G_m/H_m]/\Z\!\left(\sum_{gH_m}gH_m\right),\\
 H^1(k,T_m)=k^\times/N_{K_m/k}K_m^\times
  \longrightarrow\prod_v
      k_v^\times/N(K_m\otimes_k k_v)^\times.
\end{gathered}
\tag{A.181.10}
\]
Torus duality identifies this localization kernel with the dual
of the decomposition-group kernel in \(H^2(G_m,X_m)\),
as in \cite[Proposition 6.4]{BIR_Oki2026Norm} and the actual
inflation comparison of A.162.15.  Restriction at three is the
identity on that group, so its kernel is zero.  To spell out
the inflation step, a locally trivial lattice class restricts
over \(\Omega_m\) to a \((\Q/\Z)^{\operatorname{rank}X_m}\)
character; Chebotarev kills it.  The equality
\(H^1(G_{\Omega_m},X_m)=0\) gives inflation from the finite
group, and the same equality locally makes the full-three
restriction injective.  Thus the locally norm target \(B\)
has an actual global preimage at every finite height.
The identical argument applies to every unit target which
passes these local tests; no support restriction has entered
the lattice map A.181.10.

\paragraph{The exact one-step global-unit norm map.}
In the fixed tower, for every \(n\ge1\) one has the actual
equalities
\[
\begin{aligned}
 N_{E_{n+1}/E_n}\mathcal O_{E_{n+1}}^\times
   &=\mathcal O_{E_n}^\times\cap
                         N_{E_{n+1}/E_n}E_{n+1}^\times,\\
 N_{K_{n+1}/K_n}\mathcal O_{K_{n+1}}^\times
   &=\mathcal O_{K_n}^\times\cap
                         N_{K_{n+1}/K_n}K_{n+1}^\times.
\end{aligned}
\]
Here the intersections are taken in the displayed lower field
multiplicative group, and the maps are the actual norms.
We give the complete support correction.

Put \(L=E_{n+1}\), \(K=E_n\), and let \(\sigma\) generate
the cyclic cubic extension.  For an actual global unit
\(u=N_{L/K}z\), the divisor of \(z\) has ideal norm zero.
A nonsplit prime orbit consists of one prime, whose norm weight
is a positive integer; its divisor coefficient must therefore
be zero.  On a split three-element orbit the coefficients sum
to zero, so their vector lies in the image of \(\sigma-1\)
on the integral orbit lattice.  Taking these preimages on
the finitely many nonzero orbits produces an actual finite
fractional divisor \(D\) with
\[
                         \operatorname{div}(z)=(\sigma-1)D.
\]

A.173 for \(n=1\), and A.177 at the remaining levels,
show that the three-primary class of \(D\) is represented by
an integer divisor \(D_q\) supported on the six changed-prime
branches.  Every one of these primes is totally ramified in
\(L/K\), so \(\sigma D_q=D_q\) as an actual divisor.
The full ordinary class of \(D-D_q\) has order \(h\) prime
to three.  Choose \(y\in L^\times\) with
\(\operatorname{div}(y)=h(D-D_q)\).  The literal element
\[
                    w=\frac{z^h}{\sigma(y)/y}
\]
has zero divisor, because
\(\operatorname{div}(w)=h(\sigma-1)D_q=0\),
and has norm \(u^h\).  Choose integers \(a,b\) with
\(ah+3b=1\).  Then \(w^a u^b\) is an actual global
unit of \(L\) whose norm is exactly \(u\).
This proves the first equality.

For the second, an element \(z\in K_{n+1}^\times\)
with norm \(u\in\mathcal O_{K_n}^\times\) retains the
same norm after the quadratic scalar extension to
\(E_{n+1}/E_n\).  Apply the first equality to obtain a
unit \(w\in E_{n+1}\) with relative norm \(u\).
Put \(v=N_{E_{n+1}/K_{n+1}}w\).
Norm transitivity in the actual field square gives
\(N_{K_{n+1}/K_n}v=u^2\).
Consequently the actual global unit \(v^2/u\) has norm
\(u^{4-3}=u\), proving the second equality.
All sign and divisor identities are retained in this
quadratic transfer.

The first equality also identifies the cyclic global-unit
norm image with the units satisfying the actual local norm
conditions, by the cyclic Hasse norm theorem
\cite[Chapter VIII, Theorem 3.1, printed p.~234]{MilneCFT}.
For the pure cubic step the full decomposition subgroup at
three gives the same Hasse norm comparison through the norm
torus and its lattice, by the actual maps of A.162.12--15.
These equalities concern one specified step.  Their displayed
proof retains its lower target throughout; a later norm
condition on the chosen upper unit is an additional equation.


\emph{The exact inverse-divisor comparison above a unit target.}
Consider an actual or completed norm-compatible family with
first coordinate \(B\).  At each place of \(k\), its normalized
valuation vectors are compatible under
\[
       [\mathfrak P']\longmapsto
           f(\mathfrak P'/\mathfrak P)[\mathfrak P].
\tag{A.181.11}
\]
Their level-one coordinate is zero, since \(B\) is a global
unit.  At three and at both branches of eleven and fifty-nine,
each level is a single totally ramified field and the weight
in A.181.11 is one.  Every coordinate there is therefore zero.

At any other place with residue size \(Q\), the local algebras
are unramified.  After compatible root labels, Frobenius acts on
the root torsor \(\Z_3\) by \(F(x)=Qx+c\).
Its finite-level cycles are the primes, and the norm weight
above a cycle is the ratio of the two cycle lengths.  Every
finite orbit in \(\Z_3\) is a fixed point, because
\[
 F^h(x)-x=(1+Q+\cdots+Q^{h-1})\bigl((Q-1)x+c\bigr),
\tag{A.181.12}
\]
and the first factor is a nonzero integer.  There is at most
one fixed point.  Above any old cycle not containing it, the
three-parts of all upper cycle lengths tend uniformly to
infinity.  Otherwise periods bounded by one fixed
\(2\cdot3^r\), followed by compactness of its inverse
image in \(\Z_3\), would produce a finite orbit there.
The prime-to-three part of the periods divides two, as it
does in every finite affine group over \(\Z/3^s\).

The norm formula A.181.11 now makes each old coefficient off
the fixed branch divisible by every power of three.  The sum
over all upper cycles retains that divisibility, so moving
upper support does not change the conclusion.  Such a
coefficient is zero in both \(\Z\) and \(\Z_3\).
The only possible surviving branch has residue degree one,
and its constant coefficient is its bottom norm coordinate,
which is zero as well.  Hence every finite valuation of
every family member vanishes.

For completions this statement implies exactly membership in
the completed global units.  At a fixed level let
\(\mathcal D_m=\bigoplus_{\mathfrak P}\Z[\mathfrak P]\)
and let \(\mathcal P_m\) be the image of the principal-divisor
map.  The sequence
\[
 1\longrightarrow\mathcal U_m\longrightarrow K_m^\times
      \longrightarrow\mathcal P_m\longrightarrow0
\tag{A.181.13}
\]
splits as abelian groups, since \(\mathcal P_m\) is free.
Its pro-three completion is exact.  The completed inclusion
\(\widehat{\mathcal P}_{m,3}\to
\widehat{\mathcal D}_{m,3}\) is injective: choose an
integer \(h>0\) killing the finite class group
\(\mathcal D_m/\mathcal P_m\).  Multiplication by \(h\)
defines a map \(\mathcal D_m\to\mathcal P_m\), and its
composition with the inclusion is multiplication by \(h\).
The completion of the free group \(\mathcal P_m\) is
torsion-free, so the same composition proves injectivity.
The completion of \(\mathcal D_m\) embeds in its product
of valuation coordinates as the cofinally convergent prime
lattice.  All these coordinates were proved zero.  Thus
A.181.13 puts the original completed class in
\(\widehat{\mathcal U}_{m,3}\), uniquely and compatibly
with every norm.

We have proved literal inverse-fibre equalities with the
corresponding global-unit fibres.  Their level-three projection
would make \(U^3\) a completed global-unit norm from
\(K_3\), contrary to A.181.2.  This proves both empty
fibres in A.181.4.

\emph{The evaluated derived-limit classes.}
Define the two separate systems
\[
\begin{aligned}
 Z_m^{\rm act}&=\ker(N_{K_m/k}:K_m^\times\to k^\times),\\
 Z_m^{\rm c}&=\ker(N_{K_m/k}:\widehat{K_m^\times}_3
                                      \to\widehat{k^\times}_3).
\end{aligned}
\tag{A.181.14}
\]
Choose actual elements \(z_m\) of norm \(B\), whose existence
was proved above, and take \(z_1=B\).  The actual defects
\[
 d_m=z_m/N_{K_{m+1}/K_m}z_{m+1}\in Z_m^{\rm act}
\tag{A.181.15}
\]
define a class in \(\varprojlim^1 Z_m^{\rm act}\).
Changing the choices multiplies this sequence by
\((v_m/Nv_{m+1})_m\), so the class is independent of them.
Its image under completion lies in
\(\varprojlim^1 Z_m^{\rm c}\).  Vanishing of either
class is exactly the possibility of correcting its chosen
preimages to the corresponding norm-compatible family.
A.181.4 proves that both classes are nonzero.  The order-three
calculation in A.181.3 concerns the finite-height support
quotient; no order for the classes in A.181.15 is needed here.

\emph{The actual one-variable quotient and the continuous class.}
Let \(\Pi\) be the geometric free pro-three group on the
standard tangential generators \(x,y\) for
\(\mathbf P^1_k\setminus\{0,1,\infty\}\).
Put \(N=\ker(\Pi\to\Z_3(1))\), with \(x\mapsto1\)
and \(y\mapsto0\), and close all commutator subgroups.
The actual quotient and the map from the full metabelian
group are
\[
 P=\Pi/[N,N]=(\Lambda y)\rtimes\Z_3x,
 \quad\Lambda=\Z_3[[T]],\ X=1+T,
 \quad\ker(\Pi/\Pi''\to P)=(Y-1)D.
\tag{A.181.16}
\]
The last map is the Fox specialization of A.77.5,
\(h(X,Y)[x,y]\mapsto h(X,1)(X-1)y\).
The tangential arithmetic action sends
\(x\mapsto x^{\chi(g)}\) and
\(y\mapsto f_g^{-1}y^{\chi(g)}f_g\), with
\(f_g\in\Pi'\subset N\).  The latter conjugation
vanishes in \(N/[N,N]\), so the original arithmetic
action on the displayed quotient is retained.

Choose continuous Kummer cocycles \(a,b\) for \(A,B\),
and write \(\vartheta_g(X)=X^{\chi(g)}\).
Let \(V_A\) be \(\Lambda\) with action
\(g v=\chi(g)X^{a(g)}\vartheta_g(v)\).
The exact coefficient sequence is
\[
\begin{gathered}
 0\longrightarrow D_A\xrightarrow{\ T\ }V_A
       \xrightarrow{\epsilon}\Z_3(1)\longrightarrow0,\\
 D_A=\Lambda,\qquad
 gd=\chi(g)X^{a(g)}
       \frac{X^{\chi(g)}-1}{X-1}\vartheta_g(d).
\end{gathered}
\tag{A.181.17}
\]
The displayed action proves equivariance of multiplication
by \(T\).  A lift of the given pair to the arithmetic
extension with geometric group \(P\) is exactly a
continuous one-cocycle in \(V_A\) with augmentation
class \([b]\).  A discrepancy by an augmentation
coboundary can be removed by lifting its zero-cochain,
so its obstruction is precisely A.181.5.

For \(r\ge1\) put
\[
 f_r=X^{3^r}-1,\quad L_r=f_r/T,\quad
 V_r=V_A/f_rV_A,\quad
 J_r=\ker(V_r\xrightarrow{\epsilon}\Z_3(1)).
\tag{A.181.18}
\]
Multiplication by \(T\) identifies
\(D_A/L_rD_A=J_r\).  The finite root-torsor module
is literally
\(V_r=\Z_3[\operatorname{Roots}(Z^{3^r}-A)](1)\).
Shapiro and Kummer theory give
\[
 H^1(G_k,V_r)=\widehat{K_{r+1}^\times}_3,
 \quad \epsilon_*=N_{K_{r+1}/k},\quad
 (V_{r+1}\to V_r)_*=N_{K_{r+2}/K_{r+1}}.
\tag{A.181.19}
\]
The last root map is the literal cube map.

The finite free coefficient transitions admit continuous
\(\Z_3\)-linear sections by choosing one lift of each
root.  The \(J_r\) transitions also split as underlying
finite free modules.  Continuous cochains are consequently
degreewise surjective in both systems.
The nested compact cosets of the ideals \(f_r\Lambda\)
and \(L_r\Lambda\), whose intersection is zero, give
\(V_A=\varprojlim V_r\) and \(D_A=\varprojlim J_r\)
with their actual compact topologies.
Every \(H^0(G_k,V_r)\) is zero, by Shapiro and the
absence of an infinite compatible three-power root-of-unity
system in a number field.  Thus the continuous cochain
inverse-limit sequence gives
\[
 H^1(G_k,V_A)=\varprojlim_r
                 \widehat{K_{r+1}^\times}_3.
\tag{A.181.20}
\]
The same vanishing for \(H^0(G_k,\Z_3(1))\) identifies
\(H^1(G_k,J_r)=Z_{r+1}^{\rm c}\) and gives the injection
\[
 0\longrightarrow\varprojlim{}^1_r Z_{r+1}^{\rm c}
      \xrightarrow{\iota}H^2_{\rm cts}(G_k,D_A)
      \longrightarrow\varprojlim_rH^2_{\rm cts}(G_k,J_r)
      \longrightarrow0.
\tag{A.181.21}
\]
Omitting the initial trivial norm kernel does not change
the derived limit in A.181.14.

The sign is fixed by actual cochains.  Choose cocycles
\(h_r\in Z^1(G_k,V_r)\) for the chosen norm preimages,
and adjust them by a coboundary so that
\(\epsilon h_r=b\) literally.  If \(w\) is a continuous
lift of \(b\) to \(V_A\), the cochains
\(h_r-w_r\) lie in \(C^1(G_k,J_r)\).  With
\(\Delta(v)_r=v_r-\pi_r(v_{r+1})\), they satisfy
\(\Delta(h_r-w_r)=h_r-\pi_rh_{r+1}\) and
\(\delta(h_r-w_r)=-\delta w_r\).  Therefore
\[
                 \iota([d_m])=-\mathfrak o\ne0.
\tag{A.181.22}
\]
Every individual root-torsor image is zero by its actual
finite-height norm preimage.  At each finite completion of
\(k\), the same Shapiro and cochain maps hold for the
finite products of local root fields.  Their integral
Tate invariants are zero, and the compatible local unit
family already constructed maps to \(\kappa_3(B)\)
under augmentation.  Hence the restriction of
\(\mathfrak o\) is zero there.  At archimedean places
the higher cohomology with three-primary coefficients is
zero, so the same conclusion holds.

There is also an exact finite-coefficient comparison.
For any positive \(c,s\), the binomial identity gives
\[
 v_3\binom{3^r}{j}=r-v_3(j),\qquad
 L_r=\sum_{j=1}^{3^r}\binom{3^r}{j}T^{j-1}.
\tag{A.181.23}
\]
For sufficiently large \(r\), this puts
\(L_r\) in \((3^c,T^s)\).  The actual equivariant
quotient map
\(J_r\to D_A/(3^c,T^s)D_A\) therefore sends the
zero root-torsor obstruction to the image of
\(\mathfrak o\).  Every such finite-coefficient image
is zero although the continuous class in A.181.22 is
nonzero.  No ordinary unit-equation hypothesis is used
for this displayed pair.
\end{proof}

\paragraph{Exact certificates.}
\path{code/two_inert_evaluated_norm_phantom_certificate.py}
and its retained output and data verify the literal first-field
unit and residue identities, the norm-exponent and character
maps in A.181.6--7, the order-three support calculation, the
finite root-torsor norm weights and coefficient projections,
and finite instances of the integral cyclic divisor correction.
The all-height assertions and the nonzero continuous class are
proved by the explicit local, divisor, and cochain maps above;
the finite checks are not substituted for those arguments.
